\cotefile{157-1}
\scanpage{257}
\mspage{10}
% archive sheet number: 257
et enfin pour le troisième

\[
\begin{aligned}
&\operatorname{Hom}_{\mathrm{Hot}}(\unsure{\zeta|\xi}\times y,\unsure{\operatorname{Hom}(\mathrm{ex})|\operatorname{Hom}(\mathrm{ext})|\operatorname{Homext}|\operatorname{homext}}(M,N))\\
&\simeq\operatorname{Hom}_{\zeta\times y}(M,N)\simeq\operatorname{Hom}((\zeta\times y)\times M,N),
\end{aligned}
\]

donc toujours la même chose ! \deleted{\unsure{On gagne|On y gagne}} Pas de doutes, ça doit marcher.

S'il faut se donner le $\operatorname{Homext}$ \unsure{de (1)|de 1)} comme donnée supplémentaire, avec de plus un iso. trifonctoriel (12), il faut aussi se donner \unsure{des} iso. trifonctoriels (13), ``\uline{compatibles}'' avec les iso. (12) au sens du calcul précédent ! \deleted{\illegible{1}}

Par exemple, prenant $y={}$type d'homotopie du cercle $S^1$, les formules (13) se spécialisent en

\[
\operatorname{homext}(\Sigma M,N)\simeq
\operatorname{homext}(M,\Omega N)\simeq
\underbrace{\Omega\operatorname{Homext}(M,N)}_{\overset{\mathrm{def}}{=}\operatorname{Hom}_{\mathrm{Hot}}(S^1,-)}.
\tag*{(14)}
\]

Mais cette formule dans $\mathrm{Hot}$ n'est sans doute pas tout à fait la bonne, car le bon formalisme des $\Omega,\Sigma$ se fait dans les catégories \uline{pointées}, et en particulier dans la Catégorie homotopique pointée $\mathrm{Hot}_{\bullet}$.

\hfil$\backslash$\hfil

Je vais \unsure{expliciter} un exemple, en prenant le dérivateur

\[
\mathrm{Der}(A)
\]

\deleted{associé à} ``dérivé'' d'une catégorie abélienne $A$. Donc $\unsure{\mathbb D(e)|D(e)}={}$catégorie dérivée de $A$ --- Pour éviter des
